\documentclass[11pt, a4paper]{article}

\usepackage[T1]{fontenc}
\usepackage[utf8]{inputenc}
\usepackage{tgpagella}
\usepackage[spanish, es-noshorthands]{babel}
\usepackage{amsmath, amssymb, bm}
\usepackage{tabularx}
\usepackage{booktabs}
\usepackage[most]{tcolorbox}
\usepackage[margin=2.0cm]{geometry}
\usepackage{ragged2e}
\usepackage{fancyhdr}

\pagestyle{fancy}
\fancyhf{}
\setlength{\headheight}{16pt}
\lhead{\textbf{UCB Sede Tarija} | Ing. Mecatrónica}
\chead{}
\rhead{IMT-121 Circuitos Electrónicos I | W12S3 — Artefacto A4}
\lfoot{Prof: M.Sc. Ing. Bernardo Quiroga Turdera}
\rfoot{Página \thepage}
\renewcommand{\headrulewidth}{0.4pt}

\definecolor{ucbblue}{RGB}{0, 51, 102}
\definecolor{warnred}{RGB}{180, 0, 0}

\begin{document}

\begin{tcolorbox}[colback=warnred!5!white, colframe=warnred, title=\textbf{Artefacto A4: Checklist GO / NO-GO Docente (Autorización Física)}]
    \centering \textbf{Se Verifica por Estación y Montaje \textit{Antes} de Habilitar la Salida de Fuente}
\end{tcolorbox}

\noindent\textbf{Fecha:} \rule{2.5cm}{0.4pt} \quad \textbf{Estación:} \rule{2.5cm}{0.4pt} \quad \textbf{Docente / Téc:} \hrulefill \\
\textbf{Generador ID:} \rule{3.5cm}{0.4pt} \quad \textbf{Osciloscopio ID:} \rule{3.5cm}{0.4pt} \quad \textbf{Hora:} \rule{2cm}{0.4pt}

\section*{1. Lista de Control Crítica (Safety Gate)}
Si un solo ítem crítico no está verificado (Falla de asilamiento o de revisión), el grupo entero cae en \textbf{NO-GO}.

\vspace{0.2cm}
\renewcommand{\arraystretch}{1.55}
\noindent
\begin{tabularx}{\textwidth}{
  >{\RaggedRight\arraybackslash\hsize=1.4\hsize}X
  >{\centering\arraybackslash\hsize=0.4\hsize}X
  >{\centering\arraybackslash\hsize=0.4\hsize}X
  >{\RaggedRight\arraybackslash\hsize=1.8\hsize}X}
\toprule
\textbf{Control Crítico Obligatorio} & \textbf{Sí} & \textbf{No} & \textbf{Evidencia / Anotación de Campo} \\
\midrule
1. Equipos de la estación encendidos e inspeccionados & {[$\square$]} & {[$\square$]} & \\ \addlinespace
2. Aislamiento comprobado. Ninguna pinza GND cruzada & {[$\square$]} & {[$\square$]} & \\ \addlinespace
3. Masas de CH1/CH2 y Generador agrupadas en retorno & {[$\square$]} & {[$\square$]} & \\ \addlinespace
4. Offset, Voltajes pico y tipo de carga (Z) seguros & {[$\square$]} & {[$\square$]} & \\ \addlinespace
5. Resistores comprobados ($1/4\,\text{W}$ compatible con potencia) & {[$\square$]} & {[$\square$]} & \\ \addlinespace
6. C identificado como No Polarizado ($1\,\mu\text{F} \ge 25\,\text{V}$) & {[$\square$]} & {[$\square$]} & \\ \addlinespace
7. Piloto Docente previo con onda y frecuencia probada & {[$\square$]} & {[$\square$]} & \\ \addlinespace
8. Alumno designado para maniobrar el generador listo & {[$\square$]} & {[$\square$]} & \\
\bottomrule
\end{tabularx}

\section*{2. Resolución Formal (GO/NO-GO)}
\begin{tcolorbox}[colback=gray!5, colframe=black, title=\textbf{Toma de Decisión Documentada}]
\textbf{Dictamen Operativo (Marcar uno):} \\
{[$\square$]} \textbf{GO RC Autorizado} \hspace{1cm} {[$\square$]} \textbf{NO-GO RC} (Bloqueo de Energía) \\
{[$\square$]} \textbf{GO RLC Autorizado} \hspace{0.8cm} {[$\square$]} \textbf{NO-GO RLC} (Bloqueo de Energía)

\vspace{0.4cm}
\textbf{Justificación o Plan de Acción Inmediato (Para casos de falla):} \par\vspace{0.2cm}
\rule{\textwidth}{0.4pt} \par\vspace{0.4cm}
\rule{\textwidth}{0.4pt} \par\vspace{0.4cm}

\textbf{Contingencia Obligatoria NO-GO:} El grupo deberá mantener su estación apagada. Se limitarán a la inspección puramente desenergizada (medir resistencias y $R_{\text{DCR}}$ en frío), calcular tolerancias y leer u operar curvas en LTspice. Bajo ningún concepto deben rellenar sus tablas con inventos numéricos "ideales" fingiendo haber utilizado el osciloscopio. \par\vspace{0.3cm}

\textbf{Firma/Constancia Autorizante:} \rule{6cm}{0.4pt}
\end{tcolorbox}

\section*{3. Matriz de Observación Formativa Propuesta}
Uso del docente sin ponderación final asignada (útil para auditoría de competencias del $20\%$):\par\vspace{0.2cm}
\begin{itemize}
    \item \textbf{Predice y calcula antes de conectar:} [$\square$] Logrado \quad [$\square$] Parcial \quad [$\square$] No obs.
    \item \textbf{Identifica sondas, retornos y peligro de corto:} [$\square$] Logrado \quad [$\square$] Parcial \quad [$\square$] No obs.
    \item \textbf{Interpreta u obtiene captura auténtica:} [$\square$] Logrado \quad [$\square$] Parcial \quad [$\square$] No obs.
\end{itemize}

\end{document}
